\documentclass[a4,11pt]{wmecai}

\usepackage[brazil]{babel}      % para texto em Português
%\usepackage[english]{babel}    % para texto em Inglês

%\usepackage[latin1]{inputenc}   % para acentuação em Português OU
\usepackage[utf8]{inputenc}      % para acentuação em Português com o uso do Unicode,
                                 % mude a codificação desse padrão para utf-8

%%%%%%%%%%%%%%%%%%%%%%%%%%%%%%%%%%%%%%%%%%%%%%%%%%%%%%%%%%%%%%
%% POR FAVOR, NÃO FAÇA MUDANÇAS NESSE PADRÃO QUE ACARRETEM  EM
%% ALTERAÇÃO NA FORMATAÇÃO FINAL DO TEXTO
%%%%%%%%%%%%%%%%%%%%%%%%%%%%%%%%%%%%%%%%%%%%%%%%%%%%%%%%%%%%%%

\usepackage{graphicx}
\usepackage{subfigure}
\usepackage[centertags]{amsmath}
\usepackage{indentfirst,amsmath,amsfonts,amssymb,amsthm,newlfont}
\usepackage{longtable}
\usepackage{cite}
\usepackage[usenames]{color}
\usepackage{booktabs}


\begin{document}

%********************************************************

\title{Instruções para Submissão de Trabalhos no V WMECAI}

\author{
{\large Nome da Autora}\thanks{autora1@email}, ~{\large Nome do Autor}\thanks{autor2@email}\\
{\small ICMC-USP} \\
{\large Nome do Outro Autor }\thanks{autor3@email} \\
{\small Instituição do outro autor } \\
}

\criartitulo


%********************************************************
\section{ Introdução e Regras de Submissão}
%

No \textit{V Workshop de Matemática, Estatística e Computação Aplicadas à Indústria} (V WMECAI), os trabalhos devem ser submetidos {\bf na forma de artigos curtos (resumos expandidos)}, com {\bf 2 a 4 páginas, incluindo as referências}, de acordo com {\bf o padrão deste documento}.
Os trabalhos submetidos que não respeitem a formatação final do texto produzida por este padrão serão {\bf rejeitados} pelo Comitê Científico do evento, sem análise do mérito. Os autores podem submeter seus trabalhos em português ou inglês.

Solicitamos que não sejam alteradas as configurações de tamanho e fonte pré-deter\-mi\-na\-das deste padrão. Somente serão publicados nos anais os trabalhos {\bf aceitos e que tenham sido apresentados no evento}.

\section{Equações, Tabelas, Figuras e Referências}

As equações serão numeradas sequencialmente no texto, com a numeração à direita, usando o
comando  \verb!\label{nome-da-equacao}! para identificá-las e o comando \\
\verb!\eqref{nome-da-equacao}!, sempre que necessário
mencioná-las no texto. Por exemplo, a Equação \eqref{Calor}:
%
\begin{equation}
\frac{\partial u}{\partial t}-\Delta u = f, \quad \quad  \mathrm{em} \quad \Omega, \label{Calor}
\end{equation}
%
foi gerada pelos comandos:
\begin{verbatim}
\begin{equation}
\frac{\partial u}{\partial t}-\Delta u = f,\quad \quad  \mathrm{em}
\quad \Omega. \label{Calor}
\end{equation}
\end{verbatim}

Os autores podem introduzir figuras (veja Figura \ref{fig:exemplo}) e tabelas (veja Tabela \ref{tab:exemplo}) em seus textos. As referências bibliográficas devem seguir
o padrão adotado naquele arquivo: em ordem alfabética pelo sobrenome do primeiro autor e dos demais,
se necessário, usando-se ordem cronológica para trabalhos de um mesmo autor.
Em todas as publicações, o ano deve vir ao final, como
em \cite{Boldrini,Diniz1,Jafelice}. Exceção feita se houver  DOI, neste caso, o mesmo deve ser adicionado
após o ano de publicação, como em \cite{Cuminato,Santos}.
As referências são citadas no texto via o comando \verb!\cite{nome-da-referencia}!.
A bibliografia abaixo foi construída no ambiente \verb!\thebibliography{00}! (compila-se duas vezes o LaTeX),
com as referências  declaradas  via o comando \verb!\bibitem{nome-da-referencia}!. 

\begin{figure}[!t]
\centering
\includegraphics[width=0.5\columnwidth]{cemeai.jpg}
\caption{Exemplo de uma figura.}
\label{fig:exemplo}
\end{figure}

\begin{table}[h]
\caption{Exemplo de uma tabela.}
\centering
\begin{tabular}{l|ccc}
Dados & coluna1 & coluna2 & coluna3\\
\hline
\hline
linha1 & 110 & 220 & 330\\
\hline
linha2 & 440 & 55 & 660
\end{tabular}
\label{tab:exemplo}
\end{table}


\section{Conclusões (opcional)}
Resumir, em linhas gerais, as principais conclusões do trabalho.

\section*{Agradecimentos (opcional)}
Apresentar os agradecimentos às pessoas e instituições pertinentes, se houver espaço.

\begin{thebibliography}{00}

\bibitem{Boldrini}
J. L. Boldrini, S. I. R. Costa, V. R. Ribeiro, and H. G. Wetzler. {\it Álgebra Linear e Aplicações, 3a. edição}. Harbra, São Paulo, 1984.

\bibitem{Cuminato}
J. A. Cuminato and V. Ruas. Unification of distance inequalities for linear variational problems,
{\it Comp. Appl. Math.}, 2014. DOI: 10.1007/s40314-014-0163-6.

\bibitem{Diniz1}
G. L. Diniz, A mudança no habitat de populações de peixes: de rio a represa -- o modelo matemático, Dissertação de Mestrado em Matemática Aplicada, Unicamp, 1994.

\bibitem{Jafelice} R. M. Jafelice, L. C. Barros and R. C. Bassanezi. Study of the dynamics of HIV under
treatment considering fuzzy delay, {\it Comp. Appl. Math.}, 33:45--61, 2014.

\bibitem{Santos} I. L. D. Santos e G. N. Silva. Uma classe de problemas de controle ótimo
em escalas temporais, {\it Proceeding Series of the Brazilian Society of Computational and
Applied Mathematics}, volume 1, 2013. DOI: 10.5540/03.2013.001.01.0177.

\end{thebibliography}
\end{document}
